%!TEX root = mainQBN.tex	
\section{Quantum Factors}
\label{sec:semantic}

A \QBN\ is a DAG with nodes \emph{a set $\bX$ of \rvs} and a \emph{set $\bQ$ of quantum systems}.
We associate to each node a \emph{function}---called quantum factor---from $\Val{\bX}$ to \emph{positive operators} on the relevant Hilbert space.
We then interpret in the same way \QBN s, as well as any sub-component.
In this section we define quantum factors (\Qfactor s).
In \Cref{sec:qbn}, we reformulate the definition of \QBN s in this setting, and explain the equivalence with~\cite{GBN2014}.
Let us illustrate the idea with some examples.

\begin{example}
Consider \Cref{ex:preparation}, the preparation of a single-qubit quantum system $\qreg$ depending on a classical \rv\ $X$ (\eg\ a fair coin).
This is described by the following DAG:
\begin{tikzpicture}[scale=0.7,baseline=(b)]
\begin{genscope}
	\coordinate (b) at (0,-.2);
	\node[minimum size=5mm] (X1) at (0,0) {$X$};
	\node[minimum size=5mm] (Q1) at (1.4,0) {$\qreg$};
	\path (X1) edge (Q1);
\end{genscope}
\end{tikzpicture}.
The interpretation of this system is a function $\phi: \Val{X} \to \Pos{\HH_Q}$ giving for each $x \in \Val{X}$ a positive matrix $\M_{x}$ on $\HH_\qreg$.
If we are interested only in the state of the qubit, we can sum-out $X$ to obtain the density matrix $\Pr(x^t) \M_{x^t}+ \Pr(x^f) \M_{x^f}$, as in \Cref{ex:preparation}.

As another example, the preparation of a system $\qreg$ depending on two classical \rvs\ $X$ and $Y$ is a function from $\Val{X,Y}$ to $\Pos{\HH_\qreg}$, yielding four positive matrices $\M_{xy}\in\Pos{\HH_\qreg}$.
\end{example}

\begin{example}[Alice \& Bob]
Consider the DAG in \Cref{fig:Bell}, with four classical variables $A$, $B$, $X$ and $Y$, and one quantum node $\qreg$ producing two entangled qubits $q_1$ and $q_2$ in the state space $\HH_{q_1}\tens\HH_{q_2}$.
We associate to each node a function as follows.
\begin{itemize}
\item
To $X$ a function $\phi^X: \setofrv{X}\to\Pos{\HH_{\emptyset}}=\nnReal$ (recall \Cref{rem:pos}).
Similarly for $Y$.
\item
To $\qreg$ a positive matrix in $\Pos{\HH_{q_1}\tens\HH_{q_2}}$.
\item
To $A$ a function $\phi^A:\setofrv{X}\times\setofrv{A} \to \Pos{\HH_{q_1}}$.
\end{itemize}
\end{example}

% Intuitively, these functions generalize both the notion of factors in the theory of Bayesian networks (\ie\ functions from values of random variables to non-negative reals) and quantum instruments in quantum theory (as we will see in \cref{sec:qbn}).

\subparagraph*{Quantum registers.}

For simplicity's sake, from now on we assume quantum systems to be composed of qubits.
% All our results extend straightforwardly to the general case of classical variables of any finite arity and quantum variables of any Hilbert space.
Assuming a countable set $\mathbb{Q}$ of qubits $q_0,q_1,\dots$ with associated 2-dimensional Hilbert spaces $\HH_{q_0},\HH_{q_1},\dots$, a \textdef{quantum register} is a \emph{finite} set of qubits $\qreg \subseteq \mathbb{Q}$.
Given a quantum register $\qreg$ of $n$ qubits, its associated Hilbert space is $\HH_{\qreg} = \bigotimes_{q \in \qreg} \HH_{q} \cong \C^{2^n}$.
We write $\qreg \uplus \qreg'$ the \emph{disjoint union} of the registers $\qreg$ and $\qreg'$, only defined when $\qreg \cap \qreg' = \emptyset$.

\subparagraph*{Quantum Factors.}

Please compare (\ref{eq:qfactor}) below with (\ref{eq:factor}) the definition of a factor in \Cref{sec:BN}.

\begin{definition}[Quantum Factor (\qfactor)]
Let $\bX$ be a set of random variables and $\qreg$ a quantum register.
A \definitive{quantum factor} (shortened as \definitive{\qfactor}) $\phi$ over $(\bX,\qreg)$ is a function
\begin{equation}\label{eq:qfactor}
\someqfact: \setofrv{\bX} \to \Pos{\HH_\qreg}
\end{equation}
mapping each tuple of values $\bx$ to a \emph{positive} operator in $\L(\HH_\qreg)$.
The \definitive{scope} of $\someqfact$ is $\bX\cup\qreg$.
\end{definition}

\begin{remark}[Pure classical and pure quantum]
The two limit cases yield familiar notions.
\begin{itemize}
\item
If $\qreg$ is empty, then positive operators in $\L(\C^1)$ are exactly elements of $\nnReal$, yielding the standard definition of a factor over a set of random variables.
\item
If $\bX$ is empty, then a \qfactor\ simply corresponds to \emph{a positive operator} in $\Pos{\HH_\qreg}$.
\end{itemize}
\end{remark}

The trivial \qfactor\ over $(\emptyset,\emptyset)$ simply corresponds to a \emph{non-negative} real.
We denote by $\ftone$ the trivial \qfactor\ corresponding to $1$.

\subsection{Product and Sum of Quantum Factors}

We show that \qfactor s admit and are \emph{closed} under \emph{product} $\odot$ and \emph{summing-out} $\sum$ operations.
Remarkably, \qfactor s encompasses both tensor networks and factors from Bayesian networks:
\begin{itemize}
\item
For \qfactor s $\someqfact$ and $\someqfact'$ over classical variables only, $\someqfact \compfact \someqfact'$ is the product of factors and $\sum_X \someqfact$ is summing-out (\aka\ \emph{marginalization}).
\item
For \qfactor s $\someqfact$ and $\someqfact'$ over quantum registers only, $\someqfact \compfact \someqfact'$ is the product of tensor networks %(\aka\ contraction) 
and $\sum_\qreg \someqfact$ is tracing-out (\ie\ taking the \emph{partial trace} over $\HH_\qreg$).
\end{itemize}

\subparagraph*{Sum-out for \qfactor s.}
\emph{Summing out a variable} from a \qfactor\ corresponds to \emph{marginalization} for classical variables, and to the \emph{partial trace} for quantum systems.
\begin{definition}[Sum-out]
\label{def:sum}
Let $\someqfact$ be a \qfactor\ over $(\bX,\qreg)$. 
\begin{itemize}
\item
For $Z \in \bX$, the \definitive{sum $\sum_{Z} \someqfact $} is the \qfactor\ over $(\bX \difset\{Z\},\qreg )$ defined by:
\begin{equation*}
\left(\sum_Z \someqfact\right) (\by) \defeq \sum_{z\in\setofrv{Z}} \someqfact (\by,z)
\qquad
\text{for $\by \in \setofrv{\bX \difset\{Z\}}$}
\end{equation*}
\item
For $\qreg' \subseteq \qreg $, the \definitive{sum $\sum_{\qreg'} \someqfact $} is the \qfactor\ over $(\bX ,\qreg \difset\qreg')$ defined by:
\begin{equation*}
\left(\sum_{\qreg'} \someqfact\right) (\bx) \defeq \ptr{\HH_{\qreg'}}\left(\someqfact (\bx)\right)
\qquad
\text{for $\bx \in \setofrv{\bX}$}
\end{equation*}
\end{itemize}
\end{definition}

The definition of the product is more delicate, because classical variables and quantum systems behave differently: classical values can be ``shared'', while quantum states are linear resources.
% Such a behavior is coherent with the formalism of tensor networks, where the variables for a contraction of two tensors is given by the symmetric difference of the variables of these tensors: \eg, when contracting a tensor $T_{ijk}$ on variables $i$, $j$ and $k$ with a tensor $V_{km}$ on variables $k$ and $m$, what we obtain is a tensor $\sum_k T_{ijk}V_{km}$ on variables $i$, $j$ and $m$.
To mathematically conciliate the two, we describe the product in terms of bases.

\subparagraph{Describing \qfactor s in terms of the canonical basis.}

For a qubit $q_i$, we write $\baseof{q_i} = (e^{00}_i,e^{01}_i,e^{10}_i,e^{11}_i)$ for the canonical orthonormal basis of $\L(\HH_{q_i}) \cong \C^{2\times 2}$, \ie\ $e^{00}_i = \ket{0}\bra{0}$, $e^{01}_i = \ket{0}\bra{1}$, $e^{10}_i = \ket{1}\bra{0}$, and $e^{11}_i = \ket{1}\bra{1}$.
Given a quantum register $\qreg = \{q_1, \dots, q_m\}$, we write $\baseof{\qreg} = \prod_{i=1}^m\baseof{q_i}$, which is an orthonormal basis for $\L(\HH_\qreg) = \L(\bigotimes_{i=1}^m \HH_{q_i})$.
A \qfactor\ $\someqfact$ over $(\bX,\qreg)$ is then associated to the following function, that we denote by $\qfactfunc{\someqfact}$:
\begin{align*}
\qfactfunc{\someqfact}: \setofrv{\bX} \times \baseof{\qreg} &\to \C\\
(x_1,\dots,x_n,e_1,\dots,e_m)&\mapsto \someqfact(x_1,\dots,x_n)_{e_1,\dots,e_m}
\end{align*}
where $\someqfact(x_1,\dots,x_n)_{e_1,\dots,e_m}$ is the scalar at position $(e_1,\dots,e_m)$ of the matrix $\someqfact(x_1,\dots,x_n)$.
Notice that $\qfactfunc{\someqfact}$ uniquely determines $\someqfact$, and vice-versa.

\begin{example}
Consider \rvs\ $X_1$ and $X_2$, and a quantum register $\{q_1,q_2\}$.
An example of \qfactor\ $\someqfact$ over $(\{X_1,X_2\},\{q_1,q_2\})$ is:
\begin{equation*}
\someqfact(x_1^t,x_2^t)
=
\someqfact(x_1^t,x_2^f)
=
\someqfact(x_1^f,x_2^t)
=
\left[\begin{smallmatrix}
	1 & 0 & 0 & 0 \\
	0 & 1 & 0 & 0 \\
	0 & 0 & 1 & 0 \\
	0 & 0 & 0 & 1 \\
\end{smallmatrix}\right]
\qquad
\someqfact(x_1^f,x_2^f)
=
\left[\begin{smallmatrix}
	2 & 0 & 0 & 1 \\
	0 & 2 & 0 & 0 \\
	0 & 0 & 2 & -1 \\
	1 & 0 & -1 & 2 \\
\end{smallmatrix}\right]
\end{equation*}
%or
%$\left[\begin{smallmatrix}
%	2 & -1 & 0 & 0 \\
%	-1 & 2 & -1 & 0 \\
%	0 & -1 & 2 & i \\
%	1 & 0 & -i & 2 \\
%\end{smallmatrix}\right]$
%to have a non-real example
We also see $\someqfact$ as a function $\qfactfunc{\someqfact}:\setofrv{X_1}\times\setofrv{X_2}\times\baseof{q_1}\times\baseof{q_2} \to \C$ sending $(x_1,x_2,e_1,e_2)$ to the scalar corresponding to $e_1\tens e_2$.
For instance:
\begin{itemize}
\item
$\qfactfunc{\someqfact}(x_1^t,x_2^t,e^{00}_1,e^{00}_2) = \hphantom{-}1$,
the entry of $\someqfact(x_1^t,x_2^t)$ corresponding to $e^{00}_1 \tens e^{00}_2$;
\item
$\qfactfunc{\someqfact}(x_1^f,x_2^f,e^{11}_1,e^{01}_2) = -1$,
the entry of $\someqfact(x_1^f,x_2^f)$ corresponding to $e^{11}_1 \tens e^{01}_2$.
\end{itemize}
\end{example}

\subparagraph*{Product of \qfactor s.}
We now define the product.
It shares the values of classical variables (exactly as happens in the factors product, as discussed in \Cref{sec:BN}), while on quantum registers it behaves as the contraction of tensor networks. 
Below, we use $\qfactfunc{\someqfact}$ instead of $\someqfact$; the symmetric difference $S \symdif R \defeq (S \cup R) \difset (S \cap R)$ replaces the union in the quantum case.

\begin{definition}[Product of \qfactor s]
\label{def:prod}
Let $\someqfact_1$ and $\someqfact_2$ be \qfactor s respectively over $(\bX_1,\qreg_1)$ and $(\bX_2,\qreg_2)$.
The \definitive{product $\someqfact_1 \compfact \someqfact_2$} is a \qfactor\ over $(\bX_1\cup\bX_2,\qreg_1\symdif\qreg_2)$ defined as follows.
Let $\bx\in\setofrv{\bX_1\cup\bX_2}$ and $(\be_1,\be_2)\in\baseof{\qreg_1\symdif\qreg_2}=\baseof{\qreg_1\difset\qreg_2}\times\baseof{\qreg_2\difset\qreg_1}$.
Call $\bx_1$ the elements of $\bx$ belonging to $\setofrv{\bX_1}$, and similarly $\bx_2$ those in $\setofrv{\bX_2}$.
We set
\begin{equation*}
\qfactfunc{(\someqfact_1 \compfact \someqfact_2)} (\bx,\be_1,\be_2)
\defeq
\sum_{\be_3 \in \baseof{\qreg_1\cap\qreg_2}} \qfactfunc{\someqfact_1}(\bx_1,\be_1,\be_3)\ \qfactfunc{\someqfact_2}(\bx_2,\be_2,\be_3)
\end{equation*}
\end{definition}


\subparagraph*{Properties of sum and product.}
\qfactor s are \emph{closed under sum and product}.
Easy calculations show that the product is (essentially) associative, product and sum are both commutative, and---crucially---they distribute under suitable conditions, like classical factors.

\begin{restatable}{proposition}{qfactors}
\label{lem:qfactors_op}
\qfactor s are \emph{closed} under both product and sum.
Moreover:
\begin{enumerate}
\item\label{item:qfactors_op:1}
$(\someqfact_1\compfact \someqfact_2) \compfact \someqfact_3 = \someqfact_1 \compfact (\someqfact_2\compfact \someqfact_3)$ provided no qubit appears in the scope of all three \qfactor s
\item
$\someqfact_1 \compfact \someqfact_2 = \someqfact_2 \compfact \someqfact_1$
\item
$\sum_U \sum_V \someqfact_1 = \sum_V \sum_U \someqfact_1$
\item
$\sum_V(\someqfact_1 \compfact \someqfact_2) = (\sum_V \someqfact_1) \compfact \someqfact_2$ when $V$ is not in the scope of $\someqfact_2$
\end{enumerate}
\end{restatable}
\noindent
(Proof in \cref{app:section3}.)
The side condition in point \ref{item:qfactors_op:1} is always satisfied by quantum BNs.
% Also observe that $\ftone$ is the neutral element for the product.
